\section{Cavity dephasing due to Lamb-shift fluctuations}

The Hamiltonian of the FTT--cavity system is
\begin{align}
H(t) &= H_{\mathrm{static}} + H_{\mathrm{noise}}, \label{H}\\
H_{\mathrm{static}} &=
\omega_b\, b^\dagger b
+ \frac{K}{2}\, b^{\dagger 2} b^{2}
+ \omega_c\, c^\dagger c
+ g\,(b+b^\dagger)(c+c^\dagger), \label{Hstatic}\\
H_{\mathrm{noise}} &= \delta\omega_b(t)\, b^\dagger b .
\end{align}
Here, $b$ ($c$) annihilates an excitation in the FTT (cavity) mode, $K$ is the FTT anharmonicity, and $g$ is the coupling strength. Flux noise can cause both depolarization and dephasing. Treating the flux noise as small-amplitude $1/f$ noise, we neglect depolarization from GHz-frequency transitions (see Appendix~\ref{app:hamiltonian-derivation-noise} for details). The dominant effect is therefore an FTT-frequency fluctuation,
\begin{equation}
\delta\omega_b(t)=\delta\Phi(t)\,\partial\omega_b/\partial\Phi,
\end{equation}
where $\Phi$ is the static flux bias and $\delta\Phi(t)$ is its fluctuation.

In the dispersive regime
\begin{equation}
|g| \ll |\omega_c-\omega_b| \ll |\omega_c+\omega_b|,
\end{equation}
the counter-rotating terms in $H_{\mathrm{static}}$ can be neglected, giving
\begin{equation}
H_{\mathrm{static}}
\approx
\omega_b b^\dagger b + \omega_c c^\dagger c + \frac{K}{2} b^{\dagger 2} b^2
+ g\,(bc^\dagger + b^\dagger c).
\end{equation}
To remove the interaction term $g(bc^\dagger + b^\dagger c)$, we apply a second-order Schrieffer--Wolff transformation with $U_{\mathrm{disp}}=e^{D_{\mathrm{disp}}}$,
\begin{equation}
H_{\mathrm{static}}^{\mathrm{disp}}
=
e^{-D_{\mathrm{disp}}} H_{\mathrm{static}} e^{D_{\mathrm{disp}}}.
\end{equation}
The generator $D_{\mathrm{disp}}$ is chosen to eliminate the off-diagonal terms at first order,
\begin{equation}
\label{Ddisp}
D_{\mathrm{disp}}=
\sum_n \sqrt{n+1}\,
\frac{g}{\omega_b+nK-\omega_c}\,
|n+1\rangle\langle n|\,c
-\mathrm{h.c.}
\end{equation}
Expanding in the detuning $\Delta_{bc}=\omega_b-\omega_c$ yields
\begin{equation}
D_{\mathrm{disp}}
\approx
\frac{g}{\Delta_{bc}}\,b^\dagger c
-\frac{gK}{\Delta_{bc}^2}(b^\dagger b-1)\,b^\dagger c
+\mathcal{O}\!\left((nK/\Delta_{bc})^2\right).
\end{equation}

In the dispersive frame \cite{cqedreview2021}, the transformed static Hamiltonian takes the form
\begin{align}
\label{Hdisp}
H_\text{disp} =
\bar{\omega}_b\, b^\dagger b
+ \frac{K}{2}\, b^{\dagger 2} b^{2}
+ \bar{\omega}_c\, c^\dagger c
+ \chi\, b^\dagger b\, c^\dagger c,
\end{align}
where $\chi$ is the dispersive shift and $\bar{\omega}_b$ and $\bar{\omega}_c$ are the Lamb-shifted FTT and cavity frequencies,
\begin{align}
\chi &\approx 2\frac{g^2}{\Delta_{bc}^2}K,\quad
\bar{\omega}_b \approx \omega_b + \frac{g^2}{\Delta_{bc}}, \quad
\bar{\omega}_c \approx \omega_c - \frac{g^2}{\Delta_{bc}}.
\end{align}

The noise Hamiltonian $H_{\mathrm{noise}}$ primarily induces dephasing between eigenstates of $H_\text{disp}$ (see Sec.~\ref{subsec:Hnoise_transfer}). In particular, because $\bar{\omega}_c$ depends on the detuning $\Delta_{bc}$, FTT-frequency fluctuations $\delta\omega_b(t)$ produce cavity-frequency fluctuations,
\begin{align}
\delta\bar{\omega}_c(t)
= \frac{\partial \bar{\omega}_c}{\partial \omega_b}\,\delta\omega_b(t)
\approx \frac{g^2}{\Delta_{bc}^2}\,\delta\omega_b(t),
\end{align}
thereby generating cavity dephasing. This mechanism is summarized in Fig.~\ref{fig:device}(b).
